% Zusammenfassung (Abstract)
\section*{Zusammenfassung}
\addcontentsline{toc}{section}{Zusammenfassung}

Wer auf der Website des Immersive Reality Lab (IRL) gezielt nach Projekten, Geräten oder Publikationen sucht, stößt schnell an Grenzen: Die Inhalte sind über viele Unterseiten verteilt, eine Suchfunktion fehlt, und Fachbegriffe bleiben oft ohne Erklärung. Ergonomia begegnet diesem Problem mit einem anderen Ansatz: Statt die Navigation zu verbessern, ersetzt ein RAG-basierter Chatbot sie dort, wo eine gezielte Frage schneller zum Ziel führt als Klicken.

Entwickelt wurde das System entlang eines nutzerzentrierten Prozesses nach DIN EN ISO 9241-210. Zehn Interviews legten Nutzungssituation und Informationsbedürfnisse offen, ein erster Prototyp wurde mit denselben Nutzergruppen getestet und anschließend nachgebessert. Theoretisch orientierte sich das Projekt an ISO 9241-110, Nielsen, Norman und Amershi; technisch setzt Ergonomia auf lokale Verarbeitung mit Ollama und ChromaDB, ergänzt durch Streaming, Quellenangaben und einen dauerhaften Disclaimer zu möglichen Fehlern.

Die Nutzerstudie (N=10, Think-Aloud, SUS) ergab einen Mittelwert von 73{,}5, über dem üblichen Benchmark von 68. Aufgaben zu VR-Projekten und Publikationen wurden von nahezu allen Probanden erfolgreich abgeschlossen; neun der 18 Amershi-Guidelines sind vollständig erfüllt. Qualitativ hoben die Teilnehmenden vor allem die Präzision der Antworten und die Möglichkeit hervor, Informationen über Quellenlinks direkt zu verifizieren.

Offen bleiben Fragen, die innerhalb des Projektzeitraums nicht geklärt werden konnten. Die nach den Tests ergänzten Features (Vorschlagsfragen, Spracheingabe, Sessions, Theme) wurden nicht erneut evaluiert. Auch Barrierefreiheit, mobile Nutzung und der Umgang mit unbekannten Fachbegriffen stehen noch aus. Ergonomia ist damit kein abgeschlossenes Produkt, sondern ein erster, evaluierter Schritt mit dem Ergebnis, dass der gewählte Ansatz trägt.
